\begin{figure}[h]
\centering
\begin{tikzpicture}[decoration={
    markings,
    mark=at position 0.5 with {\arrow{>}}}
    ,scale=1.1] 

\tikzstyle{every path}=[draw] 
		\path
    node[
      regular polygon,
      regular polygon sides=12,
      draw,
      inner sep=1.9cm,
    ] (hexagon) {}
    %
    % Annotations
    (hexagon.corner 1) node[above] {$c_0$}
    (hexagon.corner 2) node[above] {}
    (hexagon.corner 3) node[left] {}
    (hexagon.corner 4) node[left] {$c_2$}
    (hexagon.corner 5) node[below] {}
    (hexagon.corner 6) node[left] {$c_4$}
    (hexagon.corner 7) node[below] {}
    (hexagon.corner 8) node[below] {$c_6$}
    (hexagon.corner 9) node[right] {$c_5$}
    (hexagon.corner 10) node[right] {$c_3$}
    (hexagon.corner 11) node[below] {}
    (hexagon.corner 12) node[right] {$c_1$}
;



\draw[postaction={decorate}] (hexagon.corner 12)--(hexagon.corner 4);
\draw[postaction={decorate}] (hexagon.corner 12)--(hexagon.corner 2);
\draw[postaction={decorate}] (hexagon.corner 12)--(hexagon.corner 3);


\draw[postaction={decorate}](hexagon.corner 4)--(hexagon.corner 10);
\draw[postaction={decorate}](hexagon.corner 4)--(hexagon.corner 11);


\draw[postaction={decorate}](hexagon.corner 6)--(hexagon.corner 9);
\draw[postaction={decorate}](hexagon.corner 10)--(hexagon.corner 5);
\draw[postaction={decorate}](hexagon.corner 10)--(hexagon.corner 6);

\draw[postaction={decorate}](hexagon.corner 9)--(hexagon.corner 7);

\draw[fill=yellow, nearly transparent] (hexagon.corner 12)--(hexagon.corner 4)--(hexagon.corner 3)--(hexagon.corner 2)--(hexagon.corner 1)--cycle;
\draw[fill=red, nearly transparent] (hexagon.corner 12)--(hexagon.corner 4)--(hexagon.corner 10)--(hexagon.corner 11)--cycle;

\draw[fill=green, nearly transparent] (hexagon.corner 10)-- (hexagon.corner 4)--(hexagon.corner 5)--(hexagon.corner 6)--cycle;
\draw[fill=blue, nearly transparent] (hexagon.corner 10)--(hexagon.corner 6)--(hexagon.corner 9)--cycle;   
\draw[fill=orange, nearly transparent] (hexagon.corner 6)--(hexagon.corner 9)--(hexagon.corner 8)--(hexagon.corner 7)--cycle;


\coordinate (m 1) at  ($0.38*(hexagon.corner 2)+0.38*(hexagon.corner 1)+0.24*(hexagon.corner 12)$);
\coordinate (m 2) at  ($0.28*(hexagon.corner 2)+0.28*(hexagon.corner 3)+0.44*(hexagon.corner 12)$);
\coordinate (m 3) at  ($0.24*(hexagon.corner 4)+0.24*(hexagon.corner 3)+0.52*(hexagon.corner 12)$);
\coordinate (m 4) at  ($0.52*(hexagon.corner 4)+0.24*(hexagon.corner 11)+0.24*(hexagon.corner 12)$);
\coordinate (m 5) at  ($0.56*(hexagon.corner 4)+0.22*(hexagon.corner 11)+0.22*(hexagon.corner 10)$);
\coordinate (m 6) at  ($0.27*(hexagon.corner 4)+0.27*(hexagon.corner 5)+0.46*(hexagon.corner 10)$);
\coordinate (m 7) at  ($0.23*(hexagon.corner 6)+0.23*(hexagon.corner 5)+0.44*(hexagon.corner 10)$);
\coordinate (m 8) at  ($0.26*(hexagon.corner 6)+0.26*(hexagon.corner 7)+0.48*(hexagon.corner 10)$);
\coordinate (m 9) at  ($0.25*(hexagon.corner 9)+0.5*(hexagon.corner 7)+0.25*(hexagon.corner 10)$);
\coordinate (m 10) at  ($0.36*(hexagon.corner 9)+0.24*(hexagon.corner 7)+0.4*(hexagon.corner 8)$);


\end{tikzpicture}
\quad
\begin{tikzpicture}[decoration={
    markings,
    mark=at position 0.5 with {\arrow{>}}}
    ,scale=1.1] 
\begin{scope}[xscale=1,yscale=1]
\tikzstyle{every path}=[draw] 
		\path
    node[
      regular polygon,
      regular polygon sides=12,
      draw=none,
      inner sep=1.9cm,
    ] (T) {}


    
    (T.corner 1) node[above] {}
    (T.corner 2) node[above] {$c_0$}
    (T.corner 3) node[left] {$c_1$}
    (T.corner 4) node[left] {}
    (T.corner 5) node[left] {$c_3$}
    (T.corner 6) node[right] {}
    (T.corner 7) node[below] {$c_{5}$}
    (T.corner 8) node[below] {$c_{4}$}
    (T.corner 9) node[left] {}
    (T.corner 10) node[right] {}
    (T.corner 11) node[right] {$c_2$}
    (T.corner 12) node[right] {}
    
;

\draw [-] (T.corner 1) to (T.corner 2) to (T.corner 3) to (T.corner 4) to (T.corner 5);
\draw [ ] (T.corner 1) to (T.corner 12);
\draw [ ] (T.corner 5) to (T.corner 6);
\draw [-] (T.corner 12)--(T.corner 11)--(T.corner 10);
\draw [ ] (T.corner 10) to (T.corner 9)to (T.corner 8);
\draw [-] (T.corner 6)--(T.corner 7)--(T.corner 8);


\coordinate (m 1) at  ($0.38*(T.corner 1)+0.38*(T.corner 2)+0.24*(T.corner 3)$);
\coordinate (m 2) at  ($0.28*(T.corner 1)+0.28*(T.corner 12)+0.44*(T.corner 3)$);
\coordinate (m 3) at  ($0.245*(T.corner 11)+0.245*(T.corner 12)+0.51*(T.corner 3)$);
\coordinate (m 4) at  ($0.52*(T.corner 11)+0.24*(T.corner 4)+0.24*(T.corner 3)$);
\coordinate (m 5) at  ($0.56*(T.corner 11)+0.22*(T.corner 4)+0.22*(T.corner 5)$);
\coordinate (m 6) at  ($0.27*(T.corner 11)+0.27*(T.corner 10)+0.46*(T.corner 5)$);
\coordinate (m 7) at  ($0.23*(T.corner 9)+0.23*(T.corner 10)+0.44*(T.corner 5)$);
\coordinate (m 8) at  ($0.26*(T.corner 9)+0.26*(T.corner 8)+0.48*(T.corner 5)$);
\coordinate (m 9) at  ($0.25*(T.corner 6)+0.5*(T.corner 8)+0.25*(T.corner 5)$);
\coordinate (m 10) at  ($0.4*(T.corner 6)+0.3*(T.corner 8)+0.3*(T.corner 7)$);


%\foreach \x in {1,...,10}{\draw (m \x) node [fill,circle,scale=0.2] {};}



\draw [name path = M ,opacity=0] (T.corner 1) to [out=15,in=40] (m 1) to [out=-180+40,in=55] (m 2) to [out=-180+55,in=75] (m 3) to [out=-180+75,in=80] (m 4) to [out=-180+80,in=90] (m 5) to [out=-180+90,in=110] (m 6) to [out=-180+110,in=125] (m 7) to [out=-180+125,in=125] (m 8) to [out=-180+125,in=105] (m 9) to [out=-180+105,in=80] (m 10) to [out=-180+80,in=80] (T.corner 8);%


\draw[name path = L 2] (T.corner 3)--(T.corner 11);
\draw[name path = L 1] (T.corner 3)--(T.corner 1);
\draw[] (T.corner 3)--(T.corner 12);
\draw[name path = L 3](T.corner 11)--(T.corner 5);
\draw[](T.corner 4)--(T.corner 11);
\draw[name path = L 4, ](T.corner 5)--(T.corner 8);
\draw[](T.corner 10)--(T.corner 5);
\draw[ ](T.corner 5)--(T.corner 9);
\draw[name path = L 5](T.corner 8)--(T.corner 6);


\path [name intersections={of= M and L 1,by=X 1}];
\path [name intersections={of= M and L 2,by=X 2}];
\path [name intersections={of= M and L 3,by=X 3}];
\path [name intersections={of= M and L 4,by=X 4}];
\path [name intersections={of= M and L 5,by=X 5}];








\draw[fill=yellow, nearly transparent] (T.corner 11)--(T.corner 12)--(T.corner 1)--(T.corner 2)--(T.corner 3)--cycle;
\draw[fill=red, nearly transparent] (T.corner 5)--(T.corner 4)--(T.corner 3)--(T.corner 11)--cycle;
\draw[fill=green, nearly transparent] (T.corner 11)--(T.corner 5)--(T.corner 8)--(T.corner 9)--(T.corner 10)--cycle;
\draw[fill=blue, nearly transparent] (T.corner 5)--(T.corner 8)--(T.corner 7)--(T.corner 6)--cycle;


\draw[postaction={decorate}] (T.corner 3)--(T.corner 1);
\draw[postaction={decorate}] (T.corner 3)--(T.corner 12);
\draw[postaction={decorate}] (T.corner 3)--(T.corner 11);
\draw[postaction={decorate}] (T.corner 11)--(T.corner 4);
\draw[postaction={decorate}] (T.corner 11)--(T.corner 5);
\draw[postaction={decorate}] (T.corner 5)--(T.corner 10);
\draw[postaction={decorate}] (T.corner 5)--(T.corner 9);
\draw[postaction={decorate}] (T.corner 5)--(T.corner 8);
\draw[postaction={decorate}] (T.corner 8)--(T.corner 6);



\end{scope}
\end{tikzpicture}
%
%
%

\caption{More examples of default orientation.}
\label{fig:default_spin_more}
\end{figure}
